\documentclass[11pt,a4paper]{article}
\usepackage[margin=2.2cm]{geometry}
\usepackage{fontspec}
\usepackage{kotex}
\setmainhangulfont{Malgun Gothic}
\setsanshangulfont{Malgun Gothic}
\usepackage{amsmath,amssymb}
\usepackage{booktabs}
\usepackage{multirow}
\usepackage{xcolor}
\usepackage{colortbl}
\usepackage{hyperref}
\usepackage{tcolorbox}
\usepackage{pgfplots}
\pgfplotsset{compat=1.18}
\tcbuselibrary{skins,breakable}
\usepackage{enumitem}
\usepackage{array}
\usepackage{float}

\definecolor{gpucol}{HTML}{7B2D8E}
\definecolor{cpubest}{HTML}{1F4E79}
\definecolor{gridwarn}{HTML}{C0392B}
\definecolor{newbest}{HTML}{D4AC0D}

\title{\LARGE\textbf{APSP 기법 통합 비교 보고서} \\[0.5em]
  \large 12 Methods $\times$ 6 Datasets + 스케일링 분석}
\author{STORM Project}
\date{2026년 4월 2일}

\begin{document}
\maketitle

%======================================================================
\section{실험 개요}
%======================================================================

본 보고서는 12종의 APSP 기법을 6개 그래프에서 비교하고, GPU 4기법의 대규모 스케일링을 분석한다. 이전 보고서(11종) 대비 다음이 변경되었다:
\begin{itemize}[leftmargin=1.4em]
  \item BG2 DAWN-SOVM 추가: DAWN(ICS 2024)의 frontier-driven BFS 재구현.
  \item TG1 D-STORM-DAWNiBFS 추가: DAWN + iBFS(SIGMOD 2016) 비트 압축을 D-STORM에 결합.
  \item TG1 D-STORM-cuBLAS 제거: $O(n^3)$ 시간 복잡도로 모든 조건에서 열세.
\end{itemize}

\subsection{프로토콜}
\begin{itemize}[leftmargin=1.4em]
  \item \textbf{측정}: 3회 실행, 최소값 보고. GPU 1회 warmup.
  \item \textbf{정확성}: 12종 전부 거리 행렬 원소별 비교 \textbf{PASS}.
  \item \textbf{TC3}: BC5와 GraphBLAS 초기화 충돌로 별도 프로세스 측정.
\end{itemize}

%======================================================================
\section{핵심 연산 용어 해설}
%======================================================================

\begin{tcolorbox}[colback=gray!3, colframe=gray!60!black, title=용어 해설, breakable]
\begin{description}[style=nextline, leftmargin=1.2em, labelwidth=0em, font=\bfseries]
\item[Edge-wise 연산] 행렬곱 없이 각 정점의 이웃을 직접 순회하여 도달 가능성을 갱신. I-AORM(BC3)에서 사용. GPU에서 CSR 이웃 순회(TG2, BG1)가 본질적으로 edge-wise.
\item[BLAS matmul] CPU 최적화 행렬곱 (SIMD, 캐시 블로킹). M-AORM(BC4), D-STORM-NumpyBLAS(TC2).
\item[SpMM] 희소 행렬곱. 비영 원소만 처리. D-STORM-SpMM-Cython(TC1)의 핵심.
\item[Cython fused prune] SpMM 출력의 COO 스캔, footprint 체크, 갱신을 단일 C 패스로 융합. TC1의 핵심 최적화.
\item[GrB\_vxm / GrB\_mxm] GraphBLAS의 벡터-행렬/행렬-행렬 곱. BC5(per-source BFS), TC3(D-STORM frontier 확장).
\item[CUDA CSR 직접 확장] frontier의 각 $(i,j)$를 CUDA 스레드가 $j$의 CSR 이웃을 직접 순회. SpMM의 세 가지 병목 제거. TG2.
\item[guard+CAS] non-atomic read로 방문 셀을 건너뛴 뒤, \texttt{atomicCAS}로 원자적 $0 \to 1$ 전이. TG2의 footprint 갱신.
\item[DAWN SOVM] Sparse Optimized Vector-Matrix Operation. frontier 정점만 확장하여 방문 완료 정점의 이웃 탐색을 건너뜀. BG2. Feng et al., ICS 2024.
\item[iBFS 비트 압축] 32개 소스의 frontier를 uint32 워드로 압축. 정점 $j$의 이웃을 1회 탐색하면서 bitwise OR로 32개 소스에 동시 배포. TG1에서 DAWN과 결합. Liu \& Huang, SIGMOD 2016.
\end{description}
\end{tcolorbox}

%======================================================================
\section{벤치마킹 대상 기법}
%======================================================================

\begin{table}[H]\centering
\caption{12종 APSP 기법 분류}
\label{tab:methods}
\small
\resizebox{\textwidth}{!}{%
\begin{tabular}{@{}clllll@{}}
\toprule
\textbf{ID} & \textbf{기법명} & \textbf{분류} & \textbf{핵심 연산} & \textbf{알고리즘} & \textbf{실행 위치} \\
\midrule
BC1 & NetworkX & CPU Baseline & Python BFS 순회 & per-source BFS & CPU (Python) \\
BC2 & SciPy & CPU Baseline & C BFS 순회 & per-source BFS & CPU (C) \\
BC3 & I-AORM & CPU Baseline & edge-wise 행 합산 & AORM 증분 & CPU (Python) \\
BC4 & M-AORM & CPU Baseline & dense BLAS matmul & AORM matmul & CPU (BLAS) \\
BC5 & GB-bfs & CPU Baseline & GrB\_vxm masked BFS & per-source BFS & CPU (C) \\
\midrule
TC1 & D-STORM-SpMM-Cython & D-STORM CPU & SciPy SpMM + Cython fused prune & frontier + footprint & CPU (Cython C) \\
TC2 & D-STORM-NumpyBLAS & D-STORM CPU & NumPy BLAS + Heaviside & frontier + footprint & CPU (BLAS) \\
TC3 & D-STORM-GraphBLAS & D-STORM CPU & GrB\_mxm + complement mask & frontier + footprint & CPU (C) \\
\midrule
\rowcolor{newbest!15} BG1 & GPU-PerSrc-BFS & GPU Baseline & CUDA block-per-source BFS & per-source BFS & GPU (CUDA) \\
\rowcolor{newbest!10} BG2 & DAWN-SOVM & GPU Baseline & CUDA frontier-driven BFS & per-source BFS & GPU (CUDA) \\
\midrule
TG1 & D-STORM-DAWNiBFS & D-STORM GPU & 비트 압축 frontier 공유 (DAWN+iBFS) & frontier + footprint & GPU (CUDA) \\
\rowcolor{newbest!10} TG2 & D-STORM-CUDA & D-STORM GPU & CUDA CSR 직접 확장 (guard+CAS) & frontier + footprint & GPU (CUDA) \\
\bottomrule
\end{tabular}}
\end{table}

%======================================================================
\section{실험 데이터 및 환경}
%======================================================================

\begin{table}[H]\centering
\caption{실험 그래프}
\label{tab:datasets}
\small
\begin{tabular}{@{}llrrrll@{}}
\toprule
\textbf{그래프} & \textbf{유형} & $n$ & $m$ & $d$ & \textbf{생성 방법} & \textbf{실험 목적} \\
\midrule
Facebook & 실제 소셜 & 4,039 & 176,468 & 8 & SNAP 수집 & 대규모 소세계 \\
BA-2000 & Scale-free & 2,000 & 19,950 & 5 & \texttt{barabasi\_albert(2000,5)} & 소세계 대표 \\
WS-2000 & Small-world & 2,000 & 12,000 & 8 & \texttt{watts\_strogatz(2000,6,0.3)} & 소세계 대표 \\
ER-2000 & Random & 2,000 & 19,882 & 6 & \texttt{erdos\_renyi(2000,0.005)} & 랜덤 그래프 \\
Grid-45$\times$45 & 2D Lattice & 2,025 & 7,920 & 88 & \texttt{grid\_2d\_graph(45,45)} & 고지름 반례 \\
PLC-2000 & PL cluster & 2,000 & 19,926 & 5 & \texttt{powerlaw\_cluster(2000,5,0.3)} & 클러스터 구조 \\
\bottomrule
\end{tabular}
\end{table}

\noindent\textbf{환경}: Windows 11 Pro, RTX 5080 (16GB), CUDA 12.9, CuPy 14.0.1, Python 3.14.3, NumPy 2.4.4, SciPy 1.17.1, Cython 3.2.4, SuiteSparse:GraphBLAS 10.3.1.

%======================================================================
\section{벤치마킹 결과}
%======================================================================

\subsection{Facebook --- CPU 기법 (8종)}

\begin{table}[H]\centering
\caption{Facebook ($n\!=\!4{,}039$): CPU 8종 순위}
\label{tab:fb-cpu}
\small
\begin{tabular}{@{}rllrrr@{}}
\toprule
\textbf{\#} & \textbf{ID} & \textbf{기법} & \textbf{시간 (s)} & \textbf{vs SciPy} & \textbf{분류} \\
\midrule
\rowcolor{cpubest!8} 1 & TC1 & D-STORM-SpMM-Cython & \textbf{0.984} & 2.0$\times$ & D-STORM \\
2 & TC2 & D-STORM-NumpyBLAS & 1.697 & 1.1$\times$ & D-STORM \\
3 & TC3 & D-STORM-GraphBLAS & 1.921 & 1.0$\times$ & D-STORM \\
4 & BC2 & SciPy & 1.918 & 1.0$\times$ & Baseline \\
5 & BC4 & M-AORM & 2.831 & 0.7$\times$ & Baseline \\
6 & BC5 & GB-bfs & 4.071 & 0.5$\times$ & Baseline \\
7 & BC3 & I-AORM & 8.693 & 0.2$\times$ & Baseline \\
8 & BC1 & NetworkX & 14.073 & 0.1$\times$ & Baseline \\
\bottomrule
\end{tabular}
\end{table}

\begin{figure}[H]
\centering
\begin{tikzpicture}
\begin{axis}[
  xbar, width=0.88\textwidth, height=7cm, bar width=9pt,
  xlabel={실행 시간 (초)},
  symbolic y coords={TC1,TC2,TC3,BC2,BC4,BC5,BC3,BC1},
  ytick=data,
  yticklabels={D-STORM-SpMM-Cython,D-STORM-NumpyBLAS,D-STORM-GraphBLAS,SciPy,M-AORM,GB-bfs,I-AORM,NetworkX},
  xmin=0, xmax=16,
  enlarge y limits={abs=8pt},
  nodes near coords={\pgfmathprintnumber[precision=3]{\pgfplotspointmeta}},
  nodes near coords style={font=\scriptsize, anchor=west},
  y tick label style={font=\small},
]
\addplot[fill=blue!40, draw=blue!60!black, forget plot]
  coordinates {(0.984,TC1) (1.697,TC2) (1.921,TC3)};
\addplot[fill=orange!18, draw=orange!50!black, bar shift=0pt, forget plot]
  coordinates {(1.918,BC2) (2.831,BC4) (4.071,BC5) (8.693,BC3) (14.073,BC1)};
\end{axis}
\end{tikzpicture}
\caption{Facebook APSP --- CPU 기법. TC1이 CPU 최고 (0.984s), SciPy 대비 2.0배.}
\label{fig:fb-cpu}
\end{figure}

\subsection{Facebook --- GPU 기법 (4종)}

\begin{table}[H]\centering
\caption{Facebook ($n\!=\!4{,}039$): GPU 4종 순위}
\label{tab:fb-gpu}
\small
\begin{tabular}{@{}rllrrr@{}}
\toprule
\textbf{\#} & \textbf{ID} & \textbf{기법} & \textbf{시간 (s)} & \textbf{vs SciPy} & \textbf{분류} \\
\midrule
\rowcolor{newbest!15} 1 & BG1 & GPU-PerSrc-BFS & \textbf{0.016} & 122$\times$ & Baseline \\
\rowcolor{newbest!10} 2 & BG2 & DAWN-SOVM & 0.019 & 101$\times$ & Baseline \\
\rowcolor{gpucol!8} 3 & TG2 & D-STORM-CUDA & 0.028 & 68$\times$ & D-STORM \\
4 & TG1 & D-STORM-DAWNiBFS & 0.247 & 7.8$\times$ & D-STORM \\
\bottomrule
\end{tabular}
\end{table}

\begin{figure}[H]
\centering
\begin{tikzpicture}
\begin{axis}[
  xbar, width=0.88\textwidth, height=5cm, bar width=10pt,
  xlabel={실행 시간 (초)},
  symbolic y coords={BG1,BG2,TG2,TG1},
  ytick=data,
  yticklabels={GPU-PerSrc-BFS (BG1),DAWN-SOVM (BG2),D-STORM-CUDA (TG2),D-STORM-DAWNiBFS (TG1)},
  xmin=0, xmax=0.30,
  enlarge y limits={abs=10pt},
  nodes near coords={\pgfmathprintnumber[precision=3]{\pgfplotspointmeta}},
  nodes near coords style={font=\scriptsize, anchor=west},
  y tick label style={font=\small},
]
\addplot[fill=yellow!50, draw=yellow!70!black, forget plot]
  coordinates {(0.016,BG1) (0.019,BG2)};
\addplot[fill=purple!40, draw=purple!60!black, bar shift=0pt, forget plot]
  coordinates {(0.028,TG2) (0.247,TG1)};
\end{axis}
\end{tikzpicture}
\caption{Facebook APSP --- GPU 기법. BG1/BG2/TG2가 Tier~1, TG1이 Tier~2.}
\label{fig:fb-gpu}
\end{figure}

\subsection{위상별 비교 --- CPU 기법}

\begin{table}[H]\centering
\caption{위상별 CPU 8기법 (초, \textbf{볼드}=해당 위상 CPU 최고)}
\label{tab:topo-cpu}
\scriptsize
\resizebox{\textwidth}{!}{%
\begin{tabular}{@{}clrrrrrr@{}}
\toprule
\textbf{ID} & \textbf{기법} & \textbf{Facebook} & \textbf{BA} & \textbf{WS} & \textbf{ER} & \textbf{Grid} & \textbf{PLC} \\
\midrule
\rowcolor{cpubest!8} TC1 & D-STORM-SpMM-Cython & \textbf{0.984} & \textbf{0.277} & \textbf{0.290} & \textbf{0.300} & \textbf{0.159} & \textbf{0.284} \\
TC2 & D-STORM-NumpyBLAS & 1.697 & 0.301 & 0.441 & 0.343 & \cellcolor{gridwarn!12}3.633 & 0.312 \\
TC3 & D-STORM-GraphBLAS & 1.921 & 0.472 & 0.493 & 0.483 & 0.590 & 0.471 \\
BC2 & SciPy & 1.918 & 0.309 & 0.291 & 0.314 & 0.161 & 0.313 \\
BC4 & M-AORM & 2.831 & 0.401 & 0.579 & 0.467 & \cellcolor{gridwarn!12}5.434 & 0.390 \\
BC5 & GB-bfs & 4.071 & 1.144 & 1.258 & 1.176 & 3.590 & 1.188 \\
BC3 & I-AORM & 8.693 & 0.393 & 0.512 & 0.430 & \cellcolor{gridwarn!12}4.461 & 0.397 \\
BC1 & NetworkX & 14.073 & 1.147 & 1.218 & 1.366 & 1.031 & 1.220 \\
\bottomrule
\end{tabular}}
\end{table}

\begin{figure}[H]
\centering
\begin{tikzpicture}
\begin{axis}[
  width=0.94\textwidth, height=8.5cm,
  xlabel={그래프 위상}, ylabel={실행 시간 (초, log scale)},
  ymode=log, ymin=0.08, ymax=30,
  log ticks with fixed point,
  symbolic x coords={Facebook,BA,WS,ER,Grid,PLC},
  xtick=data, x tick label style={font=\small},
  legend style={at={(0.50,0.98)}, anchor=north, font=\tiny, cells={anchor=west}, legend columns=2},
  grid=major, grid style={gray!25},
  mark size=3pt, thick,
]
% D-STORM CPU (blue tones, solid)
\addplot[color=blue!80!black, mark=square*, mark options={fill=blue!60}]
  coordinates {(Facebook,0.984)(BA,0.277)(WS,0.290)(ER,0.300)(Grid,0.159)(PLC,0.284)};
\addlegendentry{TC1 D-STORM-SpMM-Cython}
\addplot[color=blue!50!black, mark=triangle*, mark options={fill=blue!30}]
  coordinates {(Facebook,1.697)(BA,0.301)(WS,0.441)(ER,0.343)(Grid,3.633)(PLC,0.312)};
\addlegendentry{TC2 D-STORM-NumpyBLAS}
\addplot[color=blue!30!black, mark=diamond*, mark options={fill=blue!15}]
  coordinates {(Facebook,1.921)(BA,0.472)(WS,0.493)(ER,0.483)(Grid,0.590)(PLC,0.471)};
\addlegendentry{TC3 D-STORM-GraphBLAS}
% Baseline CPU (orange/red tones, dashed)
\addplot[color=red!70!black, mark=o, dashed]
  coordinates {(Facebook,1.918)(BA,0.309)(WS,0.291)(ER,0.314)(Grid,0.161)(PLC,0.313)};
\addlegendentry{BC2 SciPy}
\addplot[color=orange!80!black, mark=pentagon*, dashed, mark options={fill=orange!40}]
  coordinates {(Facebook,2.831)(BA,0.401)(WS,0.579)(ER,0.467)(Grid,5.434)(PLC,0.390)};
\addlegendentry{BC4 M-AORM}
\addplot[color=orange!50!black, mark=star, dashed, mark options={fill=orange!20}]
  coordinates {(Facebook,4.071)(BA,1.144)(WS,1.258)(ER,1.176)(Grid,3.590)(PLC,1.188)};
\addlegendentry{BC5 GB-bfs}
\addplot[color=red!40!black, mark=x, dashed, mark size=4pt]
  coordinates {(Facebook,8.693)(BA,0.393)(WS,0.512)(ER,0.430)(Grid,4.461)(PLC,0.397)};
\addlegendentry{BC3 I-AORM}
\addplot[color=gray!60!black, mark=+, dashed, mark size=4pt]
  coordinates {(Facebook,14.073)(BA,1.147)(WS,1.218)(ER,1.366)(Grid,1.031)(PLC,1.220)};
\addlegendentry{BC1 NetworkX}
\end{axis}
\end{tikzpicture}
\caption{위상별 CPU 기법 비교 (log scale). 실선=D-STORM, 점선=Baseline. TC1이 모든 위상에서 CPU 최고.}
\label{fig:topo-cpu}
\end{figure}

\subsection{위상별 비교 --- GPU 기법}

\begin{table}[H]\centering
\caption{위상별 GPU 4기법 (초, \textbf{볼드}=해당 위상 최고)}
\label{tab:topo-gpu}
\scriptsize
\begin{tabular}{@{}clrrrrrr@{}}
\toprule
\textbf{ID} & \textbf{기법} & \textbf{Facebook} & \textbf{BA} & \textbf{WS} & \textbf{ER} & \textbf{Grid} & \textbf{PLC} \\
\midrule
\rowcolor{newbest!12} BG1 & GPU-PerSrc-BFS & \textbf{0.016} & 0.006 & 0.005 & 0.006 & \textbf{0.006} & 0.003 \\
\rowcolor{newbest!8} BG2 & DAWN-SOVM & 0.019 & 0.006 & 0.005 & 0.006 & 0.006 & \textbf{0.003} \\
\rowcolor{gpucol!6} TG2 & D-STORM-CUDA & 0.028 & \textbf{0.005} & 0.008 & \textbf{0.005} & 0.016 & 0.005 \\
TG1 & D-STORM-DAWNiBFS & 0.247 & 0.019 & \textbf{0.009} & 0.012 & 0.032 & 0.016 \\
\bottomrule
\end{tabular}
\end{table}

\begin{figure}[H]
\centering
\begin{tikzpicture}
\begin{axis}[
  width=0.94\textwidth, height=7cm,
  xlabel={그래프 위상}, ylabel={실행 시간 (초, log scale)},
  ymode=log, ymin=0.002, ymax=0.4,
  log ticks with fixed point,
  symbolic x coords={Facebook,BA,WS,ER,Grid,PLC},
  xtick=data, x tick label style={font=\small},
  legend style={at={(0.98,0.98)}, anchor=north east, font=\scriptsize, cells={anchor=west}},
  grid=major, grid style={gray!25},
  mark size=3pt, thick,
]
% GPU Baseline (yellow/orange, dashed)
\addplot[color=yellow!70!black, mark=square*, dashed, mark options={fill=yellow!50}]
  coordinates {(Facebook,0.016)(BA,0.006)(WS,0.005)(ER,0.006)(Grid,0.006)(PLC,0.003)};
\addlegendentry{BG1 GPU-PerSrc-BFS}
\addplot[color=orange!70!black, mark=triangle*, dashed, mark options={fill=orange!40}]
  coordinates {(Facebook,0.019)(BA,0.006)(WS,0.005)(ER,0.006)(Grid,0.006)(PLC,0.003)};
\addlegendentry{BG2 DAWN-SOVM}
% D-STORM GPU (purple, solid)
\addplot[color=purple!70!black, mark=diamond*, mark options={fill=purple!40}]
  coordinates {(Facebook,0.028)(BA,0.005)(WS,0.008)(ER,0.005)(Grid,0.016)(PLC,0.005)};
\addlegendentry{TG2 D-STORM-CUDA}
\addplot[color=purple!40!black, mark=o, mark options={fill=purple!15}]
  coordinates {(Facebook,0.247)(BA,0.019)(WS,0.009)(ER,0.012)(Grid,0.032)(PLC,0.016)};
\addlegendentry{TG1 D-STORM-DAWNiBFS}
\end{axis}
\end{tikzpicture}
\caption{위상별 GPU 기법 비교 (log scale). 점선=Baseline, 실선=D-STORM. BG1/BG2/TG2가 소세계에서 수렴.}
\label{fig:topo-gpu}
\end{figure}

\subsection{전체 순위 --- CPU (6개 그래프 총합)}

\begin{table}[H]\centering
\caption{CPU 8기법: 6개 그래프 총합}
\label{tab:grand-cpu}
\small
\begin{tabular}{@{}rclrr@{}}
\toprule
\textbf{\#} & \textbf{ID} & \textbf{기법} & \textbf{총합 (s)} & \textbf{평균 (s)} \\
\midrule
\rowcolor{cpubest!5} 1 & TC1 & D-STORM-SpMM-Cython & \textbf{2.293} & 0.382 \\
2 & BC2 & SciPy & 3.305 & 0.551 \\
3 & TC3 & D-STORM-GraphBLAS & 4.429 & 0.738 \\
4 & TC2 & D-STORM-NumpyBLAS & 6.726 & 1.121 \\
5 & BC4 & M-AORM & 10.102 & 1.684 \\
6 & BC5 & GB-bfs & 12.428 & 2.071 \\
7 & BC3 & I-AORM & 14.885 & 2.481 \\
8 & BC1 & NetworkX & 20.054 & 3.342 \\
\bottomrule
\end{tabular}
\end{table}

\begin{figure}[H]
\centering
\begin{tikzpicture}
\begin{axis}[
  xbar, width=0.88\textwidth, height=7cm, bar width=9pt,
  xlabel={6개 그래프 총합 (초)},
  symbolic y coords={TC1,BC2,TC3,TC2,BC4,BC5,BC3,BC1},
  ytick=data,
  yticklabels={D-STORM-SpMM-Cython,SciPy,D-STORM-GraphBLAS,D-STORM-NumpyBLAS,M-AORM,GB-bfs,I-AORM,NetworkX},
  xmin=0, xmax=22,
  enlarge y limits={abs=8pt},
  nodes near coords={\pgfmathprintnumber[precision=2]{\pgfplotspointmeta}},
  nodes near coords style={font=\scriptsize, anchor=west},
  y tick label style={font=\small},
]
\addplot[fill=blue!40, draw=blue!60!black, forget plot]
  coordinates {(2.293,TC1) (4.429,TC3) (6.726,TC2)};
\addplot[fill=orange!18, draw=orange!50!black, bar shift=0pt, forget plot]
  coordinates {(3.305,BC2) (10.102,BC4) (12.428,BC5) (14.885,BC3) (20.054,BC1)};
\end{axis}
\end{tikzpicture}
\caption{CPU 8기법 총합. TC1이 CPU 내 최고 (2.293s).}
\label{fig:grand-cpu}
\end{figure}

\subsection{전체 순위 --- GPU (6개 그래프 총합)}

\begin{table}[H]\centering
\caption{GPU 4기법: 6개 그래프 총합}
\label{tab:grand-gpu}
\small
\begin{tabular}{@{}rclrr@{}}
\toprule
\textbf{\#} & \textbf{ID} & \textbf{기법} & \textbf{총합 (s)} & \textbf{평균 (s)} \\
\midrule
\rowcolor{newbest!12} 1 & BG1 & GPU-PerSrc-BFS & \textbf{0.041} & 0.007 \\
\rowcolor{newbest!8} 2 & BG2 & DAWN-SOVM & 0.045 & 0.008 \\
\rowcolor{gpucol!6} 3 & TG2 & D-STORM-CUDA & 0.068 & 0.011 \\
4 & TG1 & D-STORM-DAWNiBFS & 0.334 & 0.056 \\
\bottomrule
\end{tabular}
\end{table}

\begin{figure}[H]
\centering
\begin{tikzpicture}
\begin{axis}[
  xbar, width=0.88\textwidth, height=4.5cm, bar width=10pt,
  xlabel={6개 그래프 총합 (초)},
  symbolic y coords={BG1,BG2,TG2,TG1},
  ytick=data,
  yticklabels={GPU-PerSrc-BFS,DAWN-SOVM,D-STORM-CUDA,D-STORM-DAWNiBFS},
  xmin=0, xmax=0.40,
  enlarge y limits={abs=10pt},
  nodes near coords={\pgfmathprintnumber[precision=3]{\pgfplotspointmeta}},
  nodes near coords style={font=\scriptsize, anchor=west},
  y tick label style={font=\small},
]
\addplot[fill=yellow!50, draw=yellow!70!black, forget plot]
  coordinates {(0.041,BG1) (0.045,BG2)};
\addplot[fill=purple!40, draw=purple!60!black, bar shift=0pt, forget plot]
  coordinates {(0.068,TG2) (0.334,TG1)};
\end{axis}
\end{tikzpicture}
\caption{GPU 4기법 총합. BG1/BG2가 Tier~1, TG2가 Tier~1 근접, TG1이 Tier~2.}
\label{fig:grand-gpu}
\end{figure}

%======================================================================
\section{스케일링 분석}
%======================================================================

\subsection{GPU 4기법 스케일링 (BA graphs)}

\begin{table}[H]\centering
\caption{BA 그래프 스케일링: GPU 4기법 (초)}
\label{tab:scaling}
\small
\begin{tabular}{@{}rrrrr@{}}
\toprule
$n$ & \textbf{BG1} & \textbf{BG2} & \textbf{TG1} & \textbf{TG2} \\
\midrule
1,000 & \textbf{0.001} & 0.001 & 0.006 & 0.002 \\
2,000 & \textbf{0.003} & 0.003 & 0.016 & 0.005 \\
4,000 & 0.011 & \textbf{0.010} & 0.052 & 0.019 \\
6,000 & \textbf{0.023} & 0.025 & 0.100 & 0.062 \\
8,000 & \textbf{0.043} & 0.045 & 0.163 & 0.123 \\
10,000 & \textbf{0.067} & 0.072 & 0.238 & 0.211 \\
15,000 & \textbf{0.156} & 0.167 & 0.487 & 0.454 \\
\rowcolor{gridwarn!8} 20,000 & \textbf{0.263} & 0.281 & 0.805 & 12.784 \\
\bottomrule
\end{tabular}
\end{table}

\begin{figure}[H]
\centering
\begin{tikzpicture}
\begin{axis}[
  width=0.92\textwidth, height=7cm,
  xlabel={$n$ (정점 수)}, ylabel={실행 시간 (초, log scale)},
  xmin=800, xmax=22000,
  ymin=0.0008, ymax=20, ymode=log,
  log ticks with fixed point,
  legend style={at={(0.02,0.98)}, anchor=north west, font=\small},
  grid=major, grid style={gray!30},
  mark size=2.5pt,
]
\addplot[color=yellow!70!black, mark=square*, thick] coordinates
  {(1000,0.001)(2000,0.003)(4000,0.011)(6000,0.023)(8000,0.043)(10000,0.067)(15000,0.156)(20000,0.263)};
\addlegendentry{BG1 GPU-PerSrc-BFS}
\addplot[color=orange!80!black, mark=triangle*, thick] coordinates
  {(1000,0.001)(2000,0.003)(4000,0.010)(6000,0.025)(8000,0.045)(10000,0.072)(15000,0.167)(20000,0.281)};
\addlegendentry{BG2 DAWN-SOVM}
\addplot[color=blue!70!black, mark=diamond*, thick] coordinates
  {(1000,0.006)(2000,0.016)(4000,0.052)(6000,0.100)(8000,0.163)(10000,0.238)(15000,0.487)(20000,0.805)};
\addlegendentry{TG1 D-STORM-DAWNiBFS}
\addplot[color=red!70!black, mark=o, thick] coordinates
  {(1000,0.002)(2000,0.005)(4000,0.019)(6000,0.062)(8000,0.123)(10000,0.211)(15000,0.454)(20000,12.784)};
\addlegendentry{TG2 D-STORM-CUDA}
\end{axis}
\end{tikzpicture}
\caption{GPU 4기법 스케일링 (BA graphs, log scale). TG2는 $n\!=\!20{,}000$에서 메모리 절벽 (12.8초).}
\label{fig:scaling}
\end{figure}

\paragraph{분석.}
\begin{itemize}[leftmargin=1.4em]
  \item \textbf{TG2 메모리 절벽}: $n = 20{,}000$에서 TG2가 12.8초로 폭증. dense 버퍼 $16n^2 = 6.4$GB가 GPU 메모리를 압박.
  \item \textbf{TG1 안정 확장}: 비트 압축으로 메모리를 $4.25n^2$로 75\% 절감. $n = 20{,}000$에서 0.805초로 안정 동작.
  \item \textbf{BG1/BG2}: $4n^2$ 메모리, 단일 커널 launch로 전 구간 최고.
  \item \textbf{TG1 vs BG1}: 2.7배 격차는 D-STORM 프레임워크의 구조적 비용 (per-hop sync, atomicOr, 배치 관리).
\end{itemize}

\subsection{메모리 사용량 비교}

\begin{table}[H]\centering
\caption{GPU 기법별 메모리 사용량}
\label{tab:memory}
\small
\begin{tabular}{@{}llrr@{}}
\toprule
\textbf{기법} & \textbf{주요 버퍼} & \textbf{메모리} & $n\!=\!20{,}000$ \\
\midrule
BG1/BG2 & D[$n^2$] & $4n^2$ & 1.6 GB \\
TG1 & D[$n^2$] + batch 버퍼 & $\sim 4.25n^2$ & 1.7 GB \\
\rowcolor{gridwarn!8} TG2 & D + F + out\_row + out\_col & $16n^2$ & \textbf{6.4 GB} \\
\bottomrule
\end{tabular}
\end{table}

%======================================================================
\section{성능 분석}
%======================================================================

\subsection{성능 계층 구조}

\begin{tcolorbox}[colback=gray!5, colframe=gray!50!black, title=성능 계층 (전체 그래프 기준)]
\textbf{Tier~1} BG1/BG2/TG2 (0.003--0.03s) GPU BFS / DAWN / direct expand \\
\textbf{Tier~2} TG1 (0.009--0.25s) GPU 비트 압축 frontier 공유 D-STORM \\
\textbf{Tier~3} TC1/BC2 (0.16--1.92s) CPU SpMM+Cython / C BFS \\
\textbf{Tier~4} TC2/TC3/BC4 (0.30--5.43s) CPU dense / GraphBLAS \\
\textbf{Tier~5} BC1/BC3/BC5 (0.39--14.1s) Python BFS / edge-wise
\end{tcolorbox}

\subsection{D-STORM 최적화 경로: per-source BFS로의 수렴}

\begin{table}[H]\centering
\caption{D-STORM GPU 최적화 경로 --- Facebook ($n\!=\!4{,}039$) vs BA-20K ($n\!=\!20{,}000$)}
\label{tab:convergence}
\small
\begin{tabular}{@{}llrrrr@{}}
\toprule
 & & \multicolumn{2}{c}{\textbf{Facebook}} & \multicolumn{2}{c}{\textbf{BA-20K}} \\
\cmidrule(lr){3-4} \cmidrule(lr){5-6}
\textbf{단계} & \textbf{변경} & \textbf{시간} & \textbf{개선} & \textbf{시간} & \textbf{개선} \\
\midrule
TG2 DirectExpand & CSR 직접 확장 (guard+CAS) & 0.028s & --- & 12.784s & --- \\
TG1 DAWNiBFS & 비트 압축 frontier 공유 & 0.247s & 0.1$\times$ & \textbf{0.805s} & \textbf{15.9$\times$} \\
BG1 PerSrc-BFS & 프레임워크 전체 제거 & \textbf{0.016s} & 1.8$\times$ & \textbf{0.263s} & \textbf{48.6$\times$} \\
\bottomrule
\end{tabular}
\end{table}

\begin{figure}[H]
\centering
\begin{tikzpicture}
\begin{axis}[
  xbar, width=0.88\textwidth, height=5.5cm, bar width=10pt,
  xlabel={BA-20K 실행 시간 (초)},
  symbolic y coords={BG1,TG1,TG2},
  ytick=data,
  yticklabels={BG1 Per-source BFS,TG1 DAWNiBFS,TG2 DirectExpand},
  xmin=0, xmax=14,
  enlarge y limits={abs=12pt},
  nodes near coords={\pgfmathprintnumber[precision=2]{\pgfplotspointmeta}s},
  nodes near coords style={font=\scriptsize, anchor=west},
  y tick label style={font=\small},
]
\addplot[fill=yellow!50, draw=yellow!70!black, forget plot]
  coordinates {(0.263,BG1)};
\addplot[fill=purple!30, draw=purple!50!black, bar shift=0pt, forget plot]
  coordinates {(0.805,TG1)};
\addplot[fill=red!30, draw=red!60!black, bar shift=0pt, forget plot]
  coordinates {(12.784,TG2)};
\end{axis}
\end{tikzpicture}
\caption{D-STORM GPU 최적화 경로 (BA $n\!=\!20{,}000$). TG2는 메모리 절벽(12.8s)으로 실용 한계 도달. TG1은 비트 압축으로 15.9배 개선(0.81s). BG1은 프레임워크 제거로 48.6배(0.26s).}
\label{fig:convergence}
\end{figure}

\paragraph{분석.}
$n = 4{,}039$ (Facebook)에서는 TG2(0.028s)가 TG1(0.247s)보다 빠르다. 그러나 $n = 20{,}000$ (BA-20K)으로 확대하면 TG2가 메모리 절벽(12.8s)에 부딪혀 TG1(0.81s) 대비 15.9배 느려진다. 이는 TG2의 dense $16n^2$ 버퍼가 GPU VRAM을 초과하기 때문이다.

TG1(DAWNiBFS)은 비트 압축으로 메모리를 75\% 절감하여 대규모에서 안정 동작하나, BG1(0.26s) 대비 3.1배 느리다. 이 격차는 D-STORM 프레임워크의 구조적 비용(배치 루프, per-hop sync, atomicOr)이며 엔지니어링이 아닌 \textbf{구조적 한계}이다.

%======================================================================
\section{연구 기여}
%======================================================================

\begin{tcolorbox}[colback=green!5, colframe=green!50!black, title=연구 기여, breakable]

본 연구는 ``D-STORM이 가장 빠르다''는 주장이 아니라, \textbf{행렬 대수 APSP의 한계와 가능성을 체계적으로 규명}한 것이 핵심이다.

\begin{enumerate}[leftmargin=1.4em]
  \item \textbf{행렬 대수 APSP가 CPU에서 C BFS를 능가}: TC1(D-STORM-SpMM-Cython)이 SciPy C BFS를 6개 위상 전부에서 능가 (최대 2.0$\times$). Python 오케스트레이션의 행렬 대수 접근이 최적화된 C 루프를 이기는 비자명한 결과.

  \item \textbf{GPU 행렬 대수 APSP는 per-source BFS로 수렴}: D-STORM의 SpMM 병목을 순차 제거하면 per-source BFS에 수렴. 행렬 대수 프레임워크의 간접 비용(공유 footprint, per-hop sync, 배치 관리)이 순수 BFS 대비 불가피한 오버헤드.

  \item \textbf{12종 체계적 벤치마크}: Python BFS, C BFS, AORM (edge-wise/matmul), D-STORM (CPU 3종 + GPU 2종), GraphBLAS (BFS/D-STORM), GPU per-source BFS, DAWN을 동일 환경에서 비교. 기존 문헌에서 찾기 어려운 공정한 교차 패러다임 비교.

  \item \textbf{D-STORM 고유 기능}: 동적 간선 삽입 ($D'_{ij} = \min(D_{ij}, D_{ia}+1+D_{bj})$, $O(n^2)$ 갱신, 22--33$\times$ 가속)은 BFS로 대체 불가. 대수적 수렴 분석 프레임워크.

  \item \textbf{GPU D-STORM 확장성}: TG1(DAWNiBFS)이 iBFS 비트 압축으로 TG2의 메모리 병목($16n^2$)을 해결하여 $n = 20{,}000$+ 동작 가능. BFS 대비 2.7$\times$ 느리나, D-STORM 프레임워크를 유지하면서 대규모 동작하는 최선의 구현.
\end{enumerate}
\end{tcolorbox}

%======================================================================
\section{종합 결론}
%======================================================================

\begin{tcolorbox}[colback=blue!3, colframe=blue!50!black, title=조건부 운영 가이드]
\begin{itemize}[leftmargin=1.4em]
  \item \textbf{GPU, full APSP}: BG1 GPU-PerSrc-BFS (모든 조건 최고).
  \item \textbf{GPU, D-STORM 필요 (동적 갱신 등)}: $n < 15{,}000$ $\to$ TG2, $n \geq 15{,}000$ $\to$ TG1.
  \item \textbf{CPU 전용}: TC1 D-STORM-SpMM-Cython (모든 위상 CPU 최고, SciPy 대비 최대 2.0$\times$).
  \item \textbf{동적 간선 삽입}: D-STORM (TC1/TG1/TG2) --- BFS 대체 불가.
\end{itemize}
\end{tcolorbox}

%======================================================================
\section{향후 연구: D-STORM 동적 갱신의 응용}
%======================================================================

D-STORM의 동적 간선 삽입($D'_{ij} = \min(D_{ij}, D_{ia}+1+D_{bj})$, $O(n^2)$)은 full APSP 재계산($O(n \cdot m)$) 대비 22--33배 빠르며, per-source BFS로 대체할 수 없는 고유 기능이다. 이 기능이 실질적 가치를 갖는 응용 분야를 정리한다.

\begin{tcolorbox}[colback=gray!3, colframe=gray!60!black, title=동적 갱신 응용 계획, breakable]
\begin{description}[style=nextline, leftmargin=1.2em, labelwidth=0em, font=\bfseries]

\item[1. 소셜 네트워크 실시간 영향력 분석]
새로운 친구 관계(간선)가 추가될 때마다 전체 사용자 간 거리가 변화한다. Closeness centrality, betweenness centrality 등 거리 기반 지표가 간선 추가마다 변동하므로, 영향력 순위의 실시간 추적이 필요하다. 간선 100개 연속 추가 시 BFS 전체 재계산은 $100 \times O(n \cdot m)$이나, D-STORM 증분 갱신은 $100 \times O(n^2)$으로 1--2자릿수 빠르다.

\item[2. 교통/물류 네트워크 경로 갱신]
새 도로 개통, 노선 신설 등 네트워크 구조 변경 시 모든 지점 간 최단 경로를 즉시 반영해야 하는 실시간 경로 안내 시스템. 간선 추가가 빈번하고 전체 재계산 비용이 큰 환경에서 D-STORM의 증분 갱신이 유리하다. 특히 소세계 특성을 가진 도시 교통 네트워크(d $\leq$ 10)에서 효과적이다.

\item[3. 통신 네트워크 토폴로지 관리]
SDN(Software-Defined Networking) 환경에서 라우터/링크 추가 시 라우팅 테이블 갱신, 장애 복구 후 링크 재연결 시 경로 재계산 등에 활용 가능하다. 네트워크 규모가 수천 노드이고 변경이 빈번한 데이터센터 환경에 적합하다.

\item[4. 전염병 확산 경로 추적]
접촉 네트워크에서 새 접촉(간선)이 발생할 때마다 감염 경로 거리를 갱신하여 감염 위험도를 실시간 평가한다. 접촉 추적 시스템에서 접촉 이벤트는 연속적으로 발생하므로, 매 이벤트마다 전체 BFS를 수행하는 것은 비현실적이다.

\item[5. 금융 거래 네트워크 위험 전파 분석]
금융 기관 간 거래 관계가 신설될 때 시스템 리스크 전파 경로의 변화를 추적한다. 새 거래 관계(간선)의 추가가 전체 네트워크 연결성과 위험 전파 거리에 미치는 영향을 실시간으로 평가할 수 있다.

\end{description}

\paragraph{공통 특성.}
위 응용들은 (1) 간선이 \textbf{연속적으로 추가}되고, (2) 추가 즉시 \textbf{전체 거리 행렬 갱신}이 필요하며, (3) \textbf{소세계 구조}($d \leq 10$)를 가진다는 공통점이 있다. 이 조건에서 D-STORM의 $O(n^2)$ 증분 갱신은 $O(n \cdot m)$ 전체 재계산 대비 1--2자릿수 빠르다.

\paragraph{향후 실험 계획.}
\begin{enumerate}[leftmargin=1.4em]
  \item 실제 시계열 소셜 네트워크 데이터(예: Temporal Facebook, arXiv collaboration)에서 간선 스트림 처리 성능 측정.
  \item 간선 추가 빈도(초당 간선 수)에 따른 BFS 재계산 vs D-STORM 증분 갱신의 처리량(throughput) 비교.
  \item D-STORM 동적 갱신의 GPU 구현(TG1/TG2 기반)으로 대규모 동적 그래프 대응.
  \item 간선 \textbf{삭제}에 대한 D-STORM 확장 연구 (아래 상세).
\end{enumerate}
\end{tcolorbox}

\subsection{간선 삭제 알고리즘 설계}

간선 삽입은 거리가 감소만 하므로 $D'_{ij} = \min(D_{ij}, D_{ia}+1+D_{bj})$로 $O(n^2)$에 해결된다. 반면 간선 삭제는 거리가 증가할 수 있어, 삭제된 간선이 최단 경로에 사용되었는지 판정하고 대체 경로를 탐색해야 한다.

본 알고리즘의 핵심인 ``대체 부모 체크''는 Even--Shiloach~\cite{even1981} 알고리즘에 기반한다. Even--Shiloach는 BFS 트리를 유지하며 간선 삭제 시 레벨 증가를 자식에게 점진적으로 전파하여 단일 소스에 대해 총 $O(m \cdot n)$을 달성한다. 본 설계는 이를 D-STORM의 거리 행렬 $D$에 적용한 \textbf{단순화 변형}으로, BFS 트리 유지 없이 $D$ 값만으로 영향 소스를 식별하고, 영향받는 소스에 대해서만 BFS 행을 재계산한다.

\begin{table}[H]\centering
\caption{관련 동적 APSP 알고리즘}
\small
\begin{tabular}{@{}llll@{}}
\toprule
\textbf{알고리즘} & \textbf{핵심 아이디어} & \textbf{복잡도} & \textbf{비고} \\
\midrule
Even--Shiloach (1981) & BFS 트리 + 대체 부모 + 레벨 전파 & $O(m \cdot n)$ 총합 & 단일 소스, 감소적 \\
Ramalingam--Reps (1996) & 영향 정점만 재계산 & 영향 비례 & 삽입+삭제 \\
Demetrescu--Italiano (2004) & 행렬 기반 완전 동적 & $O(n^2 \log^3 n)$ 상각 & 가중치 지원 \\
\midrule
\textbf{본 설계} & D-STORM $D$ + 대체 부모 + 부분 BFS & $O(n \cdot d + |A| \cdot m)$ & 구현 단순, $D$ 재활용 \\
\bottomrule
\end{tabular}
\end{table}

Even--Shiloach 대비 본 설계의 차이점: (1) BFS 트리를 유지하지 않아 메모리 절약, (2) 레벨 전파 대신 영향 소스의 전체 행을 BFS로 재계산하여 구현이 단순, (3) D-STORM의 기존 거리 행렬 $D$를 직접 활용하여 별도 자료구조 불필요. 다만 레벨 전파가 없으므로 연쇄적 거리 증가 시 Even--Shiloach보다 비효율적일 수 있다.

\paragraph{핵심 아이디어: 대체 부모 체크.}
간선 $(a,b)$ 삭제 시, 소스 $s$에서 $b$까지의 최단 경로가 $(a,b)$를 사용했다면 $D(s,a)+1 = D(s,b)$이다. 이때 $b$의 이웃 중 $a$가 아닌 정점 $u$에 대해 $D(s,u) = D(s,b)-1$인 \textbf{대체 부모}가 존재하면, $(a,b)$ 없이도 같은 거리로 $b$에 도달 가능하므로 소스 $s$는 영향받지 않는다.

\begin{tcolorbox}[colback=blue!3, colframe=blue!50!black, title=간선 삭제 알고리즘]
\begin{verbatim}
DELETE_EDGE(D, A, a, b):
    A.remove_edge(a, b)
    affected = []

    for s in range(n):
        # 방향 1: s -> ... -> a -> b
        if D[s,a] + 1 == D[s,b]:
            if not any(D[s,u]==D[s,b]-1 for u in neighbors(b) if u!=a):
                affected.append(s); continue

        # 방향 2: s -> ... -> b -> a
        if D[s,b] + 1 == D[s,a]:
            if not any(D[s,u]==D[s,a]-1 for u in neighbors(a) if u!=b):
                affected.append(s)

    # 영향받는 소스만 BFS 재계산
    for s in affected:
        D[s,:] = BFS(A, s)
\end{verbatim}
\end{tcolorbox}

\paragraph{복잡도 분석.}

\begin{table}[H]\centering
\caption{간선 삭제 복잡도}
\small
\begin{tabular}{@{}llll@{}}
\toprule
\textbf{단계} & \textbf{복잡도} & \textbf{설명} \\
\midrule
영향 소스 식별 & $O(n \cdot d_{\text{avg}})$ & $n$개 소스 $\times$ 이웃 대체 부모 체크 \\
BFS 재계산 & $O(|\text{affected}| \cdot m)$ & 영향 소스만 BFS \\
\midrule
\textbf{총합} & $O(n \cdot d + |\text{affected}| \cdot m)$ & \\
\bottomrule
\end{tabular}
\end{table}

\begin{table}[H]\centering
\caption{$|\text{affected}|$의 그래프 유형별 예상}
\small
\begin{tabular}{@{}lll@{}}
\toprule
\textbf{시나리오} & $|\textbf{affected}|$ & \textbf{총 비용} \\
\midrule
소세계, 고연결 & $O(1)$ & $O(n \cdot d)$ --- 거의 공짜 \\
소세계, 평균 & $O(\sqrt{n})$ & $O(n \cdot d + \sqrt{n} \cdot m)$ \\
트리 (최악) & $O(n)$ & $O(n \cdot m)$ --- 전체 재계산과 동일 \\
\bottomrule
\end{tabular}
\end{table}

\paragraph{소세계 그래프에서의 유효성.}
소세계 네트워크는 높은 클러스터링 계수를 가지므로, 대부분의 정점이 다수의 대체 경로를 보유한다. 따라서 간선 하나의 삭제가 최단 경로를 변경하는 소스의 비율($|\text{affected}|/n$)은 극히 낮을 것으로 예상된다. 이 가설은 향후 실제 시계열 소셜 네트워크 데이터에서 $|\text{affected}|$의 분포를 측정하여 검증할 예정이다.

\paragraph{삽입+삭제 통합 프레임워크.}
삽입과 삭제를 결합하면 D-STORM은 \textbf{완전 동적(fully dynamic)} APSP 시스템으로 확장된다:
\begin{itemize}[leftmargin=1.4em]
  \item \textbf{삽입}: $O(n^2)$ 닫힌 형태 갱신 (기존).
  \item \textbf{삭제}: $O(n \cdot d + |\text{affected}| \cdot m)$ 대체 부모 체크 + 부분 BFS (제안).
  \item \textbf{조합}: 간선 스트림에서 삽입/삭제가 혼재하는 동적 네트워크에 대응.
\end{itemize}

\begin{thebibliography}{9}
\bibitem{even1981} S. Even and Y. Shiloach, ``An On-line Edge-Deletion Problem,'' \textit{Journal of the ACM}, vol.~28, no.~1, pp.~1--4, 1981.
\bibitem{ramalingam1996} G. Ramalingam and T. Reps, ``An Incremental Algorithm for a Generalization of the Shortest-Path Problem,'' \textit{Journal of Algorithms}, vol.~21, no.~2, pp.~267--305, 1996.
\bibitem{demetrescu2004} C. Demetrescu and G.~F. Italiano, ``A New Approach to Dynamic All Pairs Shortest Paths,'' \textit{Journal of the ACM}, vol.~51, no.~6, pp.~968--992, 2004.
\end{thebibliography}

\end{document}
